\section{Anexo: Configuración del Backend}
\label{sec:anexo_setup_backend}

Comandos para crear los recursos AWS necesarios \textbf{una sola vez} antes del primer despliegue: buckets S3 para código, checkpoints, datasets y estado de Terraform.

% ---------------------------------------------------------------
\subsection{Prerequisitos}

\begin{lstlisting}[language=bash, caption={Verificación de credenciales AWS.}]
# Verificar que tienes AWS CLI y credenciales configuradas
aws sts get-caller-identity
# Debe mostrar tu Account ID y User ARN
\end{lstlisting}

% ---------------------------------------------------------------
\subsection{1. Crear Buckets S3}

\subsubsection{1a. Bucket para Código (\texttt{tfm-slm-code})}
\texttt{deploy.py} archiva y sube el código aquí. EC2 lo descarga al iniciar.

\begin{lstlisting}[language=bash]
aws s3api create-bucket \
    --bucket tfm-slm-code \
    --region eu-south-2 \
    --create-bucket-configuration LocationConstraint=eu-south-2
\end{lstlisting}

\subsubsection{1b. Bucket para Checkpoints (\texttt{tfm-slm-checkpoints})}
El entrenamiento guarda checkpoints aquí tras cada época. La inferencia los descarga si no están en local.

\begin{lstlisting}[language=bash]
aws s3api create-bucket \
    --bucket tfm-slm-checkpoints \
    --region eu-south-2 \
    --create-bucket-configuration LocationConstraint=eu-south-2
\end{lstlisting}

\subsubsection{1c. Bucket para Datasets (\texttt{tfm-slm-datasets})}
El dataset procesado (\texttt{.arrow}) se sube aquí tras el primer procesado. Runs posteriores lo restauran directamente, evitando reprocesar las 346.439 muestras.

\begin{lstlisting}[language=bash]
aws s3api create-bucket \
    --bucket tfm-slm-datasets \
    --region eu-south-2 \
    --create-bucket-configuration LocationConstraint=eu-south-2
\end{lstlisting}

\subsubsection{1d. Bucket para Benchmarks (\texttt{tfm-slm-benchmarks})}
Almacena resultados de benchmarking (métricas, perplejidad, throughput) en JSON.

\begin{lstlisting}[language=bash]
aws s3api create-bucket \
    --bucket tfm-slm-benchmarks \
    --region eu-south-2 \
    --create-bucket-configuration LocationConstraint=eu-south-2
\end{lstlisting}

\subsubsection{1e. Bucket para Estado Terraform (\texttt{tfm-slm-terraform-state})}
Almacena \texttt{terraform.tfstate} de forma centralizada, permitiendo múltiples ejecuciones seguras.

\begin{lstlisting}[language=bash]
aws s3api create-bucket \
    --bucket tfm-slm-terraform-state \
    --region eu-south-2 \
    --create-bucket-configuration LocationConstraint=eu-south-2
\end{lstlisting}

% ---------------------------------------------------------------
\subsection{2. Habilitar Versionado}

Habilita versionado en el bucket de estado para recuperar versiones anteriores si algo falla.

\begin{lstlisting}[language=bash, caption={Versionado en bucket de estado Terraform.}]
aws s3api put-bucket-versioning \
    --bucket tfm-slm-terraform-state \
    --versioning-configuration Status=Enabled
\end{lstlisting}

% ---------------------------------------------------------------
\subsection{3. Validar Buckets Creados}

\begin{lstlisting}[language=bash, caption={Validación de los buckets creados en eu-south-2.}]
# Listar todos los buckets
aws s3 ls

# Verificar que todos los buckets estan en eu-south-2
aws s3api list-buckets --query 'Buckets[?contains(Name, `tfm-slm`)]'

# Verificar versionado habilitado en terraform-state
aws s3api get-bucket-versioning --bucket tfm-slm-terraform-state
\end{lstlisting}

% ---------------------------------------------------------------
\subsection{4. Flujo Completo tras el Setup}

\begin{lstlisting}[language=bash, caption={Flujo completo de despliegue tras la configuración inicial.}]
# 1. Preparar codigo y subirlo a S3 (1=train, 2=inference, 3=benchmark)
echo "1" | python3 deploy.py

# 2. Provisionar EC2 (descarga codigo, construye Docker con GPU)
cd app/terraform
terraform init       # Solo la primera vez
terraform apply -auto-approve -lock=false

# 3. Monitorizar build y entrenamiento
ssh -i "tfm-slm.pem" ec2-user@<IP> tail -f /var/log/user-data.log
\end{lstlisting}

\subsection*{Notas}
\begin{itemize}
    \item \textbf{S3 Native Locking}: Terraform usa \texttt{use\_lockfile}~=~\texttt{true} (no requiere DynamoDB).
    \item \textbf{Credenciales AWS}: Se leen automáticamente desde \texttt{\textasciitilde/.aws/credentials}.
    \item \textbf{Región}: Todos los buckets en \texttt{eu-south-2}.
    \item \textbf{Nombres únicos}: Los nombres de buckets S3 son globales en AWS. Si ya existen, usar nombres distintos en \texttt{deploy.py} y \texttt{terraform.tfvars}.
\end{itemize}
